% GENERATED FILE. Edit canonical lessons, then run npm run generate:books.
% canonical-source-hash: fnv1a64:92dd74ab302087e7
% canonical-lessons: JA-C307-tsunagaru, JA-C307-tsunagu, JA-C307-dengon, JA-C307-keitai, JA-C307-kakenaosu

\chapter{A Call in the Night}
\label{ch:ja-a-call-in-the-night}
\hlchaptermodality{307}

\begin{chapteropening}
\textbf{By the end of this chapter:} \emph{I can say whether a call got through and leave a message.}

Complete the last lesson of chapter 307: I can say whether a call got through and leave a message.
\end{chapteropening}

\section[\texorpdfstring{\ja{つながる} (tsunagaru)}{つながる (tsunagaru)}]{\ja{つながる} (tsunagaru) — to be connected; to get through}

\label{lesson:JA-C307-tsunagaru}



\begin{quote}
Before the new one: say the Japanese for surface mail, then the Japanese for to receive.
\end{quote}

\subsection*{You'll want to know: \ja{つながる}}
\textbf{\ja{つながる}} — \emph{tsunagaru} — \textquotedblleft{}to be connected; to get through\textquotedblright{}.

\begin{quote}
\textbf{\ja{やっと} \ja{でんわが} \ja{つながりました。}} — \emph{yatto denwa ga tsunagarimashita} — \textquotedblleft{}At last the call went through.\textquotedblright{}
\end{quote}

Polite \textbf{\ja{つながります}}, \emph{tsunagarimasu}; a godan verb. Past (polite): \textbf{\ja{でんわが} \ja{つながりました。}}, \emph{denwa ga tsunagarimashita}.

\begin{cousinweb}
\textbf{\ja{つながる}} (\emph{tsunagaru}) is a native Japanese word (wago), inherited rather than borrowed from Chinese; its kanji was attached later as a kun'yomi (JMdict).
\end{cousinweb}

\begin{grammarlens}[title={What you've built}]
One more word for telling what happened.
\end{grammarlens}

\subsection*{Your turn}
\begin{itemize}
\raggedright
  \item \emph{Say it:} \emph{tsunagaru}
  \item \emph{Say it:} \emph{tsunagaru}, once more
  \item \emph{Say it:} \emph{yatto denwa ga tsunagarimashita} — At last the call went through.
  \item \emph{Recall:} say the Japanese for surface mail, then the Japanese for to receive, then say \emph{tsunagaru} again
\end{itemize}

\begin{tcolorbox}[breakable,colback=teal!4,colframe=teal!35!black,title={Before you move on}]
What does \ja{つながる} mean? (\textquotedblleft{}To be connected; to get through\textquotedblright{}.) Say it once more.
\end{tcolorbox}
\section[\texorpdfstring{\ja{つなぐ} (tsunagu)}{つなぐ (tsunagu)}]{\ja{つなぐ} (tsunagu) — to join; to hold (hands)}

\label{lesson:JA-C307-tsunagu}



\begin{quote}
Before the new one: say the Japanese for to receive, then the Japanese for to be connected.
\end{quote}

\subsection*{You'll want to know: \ja{つなぐ}}
\textbf{\ja{つなぐ}} — \emph{tsunagu} — \textquotedblleft{}to join; to hold (hands)\textquotedblright{}.

\begin{quote}
\textbf{\ja{こどもと} \ja{てを} \ja{つないで} \ja{あるきました。}} — \emph{kodomo to te o tsunaide arukimashita} — \textquotedblleft{}I walked hand in hand with the child.\textquotedblright{}
\end{quote}

Polite \textbf{\ja{つなぎます}}, \emph{tsunagimasu}; a godan verb. Past (polite): \textbf{\ja{こどもが} \ja{てを} \ja{つなぎました。}}, \emph{kodomo ga te o tsunagimashita}.

\begin{cousinweb}
\textbf{\ja{つなぐ}} (\emph{tsunagu}) is a native Japanese word (wago), inherited rather than borrowed from Chinese; its kanji was attached later as a kun'yomi (JMdict). It is the transitive partner of \emph{tsunagaru}, \textquotedblleft{}to be connected\textquotedblright{}.
\end{cousinweb}

\begin{grammarlens}[title={What you've built}]
One more word for telling what happened.
\end{grammarlens}

\subsection*{Your turn}
\begin{itemize}
\raggedright
  \item \emph{Say it:} \emph{tsunagu}
  \item \emph{Say it:} \emph{tsunagu}, once more
  \item \emph{Say it:} \emph{kodomo to te o tsunaide arukimashita} — I walked hand in hand with the child.
  \item \emph{Recall:} say the Japanese for to receive, then the Japanese for to be connected, then say \emph{tsunagu} again
\end{itemize}

\begin{tcolorbox}[breakable,colback=teal!4,colframe=teal!35!black,title={Before you move on}]
What does \ja{つなぐ} mean? (\textquotedblleft{}To join; to hold (hands)\textquotedblright{}.) Say it once more.
\end{tcolorbox}
\section[\texorpdfstring{\ja{でんごん} (dengon)}{でんごん (dengon)}]{\ja{でんごん} (dengon) — a message (passed on)}

\label{lesson:JA-C307-dengon}



\begin{quote}
Before the new one: say the Japanese for to be connected, then the Japanese for to join.
\end{quote}

\subsection*{You'll want to know: \ja{でんごん}}
\textbf{\ja{でんごん}} — \emph{dengon} — \textquotedblleft{}a message (passed on)\textquotedblright{}.

\begin{quote}
\textbf{\ja{あにから} \ja{でんごんが} \ja{あります。}} — \emph{ani kara dengon ga arimasu} — \textquotedblleft{}There is a message from my brother.\textquotedblright{}
\end{quote}

\begin{cousinweb}
\textbf{\ja{でんごん}} (\emph{dengon}) is Sino-Japanese (kango): two Chinese characters read with their on'yomi, readings borrowed from Chinese. They are \emph{den}, \textquotedblleft{}transmit\textquotedblright{} (Middle Chinese \emph{drjwen}) and \emph{gon}, \textquotedblleft{}word\textquotedblright{} (\emph{ngjon}) (KANJIDIC2; Middle Chinese in Baxter's transcription).
\end{cousinweb}

\begin{grammarlens}[title={What you've built}]
One more word for telling what happened.
\end{grammarlens}

\subsection*{Your turn}
\begin{itemize}
\raggedright
  \item \emph{Say it:} \emph{dengon}
  \item \emph{Say it:} \emph{dengon}, once more
  \item \emph{Say it:} \emph{ani kara dengon ga arimasu} — There is a message from my brother.
  \item \emph{Recall:} say the Japanese for to be connected, then the Japanese for to join, then say \emph{dengon} again
\end{itemize}

\begin{tcolorbox}[breakable,colback=teal!4,colframe=teal!35!black,title={Before you move on}]
What does \ja{でんごん} mean? (\textquotedblleft{}A message (passed on)\textquotedblright{}.) Say it once more.
\end{tcolorbox}
\section[\texorpdfstring{\ja{けいたい} (keitai)}{けいたい (keitai)}]{\ja{けいたい} (keitai) — a mobile phone}

\label{lesson:JA-C307-keitai}



\begin{quote}
Before the new one: say the Japanese for to join, then the Japanese for a message (passed on).
\end{quote}

\subsection*{You'll want to know: \ja{けいたい}}
\textbf{\ja{けいたい}} — \emph{keitai} — \textquotedblleft{}a mobile phone\textquotedblright{}.

\begin{quote}
\textbf{\ja{けいたいを} \ja{うちに} \ja{わすれました。}} — \emph{keitai o uchi ni wasuremashita} — \textquotedblleft{}I left my mobile at home.\textquotedblright{}
\end{quote}

\begin{cousinweb}
\textbf{\ja{けいたい}} (\emph{keitai}) is Sino-Japanese (kango): two Chinese characters read with their on'yomi, readings borrowed from Chinese. They are \emph{kei}, \textquotedblleft{}carry (in hand)\textquotedblright{} (Middle Chinese \emph{hwej}) and \emph{tai}, \textquotedblleft{}belt\textquotedblright{} (\emph{tajH}) (KANJIDIC2; Middle Chinese in Baxter's transcription).
\end{cousinweb}

\begin{grammarlens}[title={What you've built}]
One more word for telling what happened.
\end{grammarlens}

\subsection*{Your turn}
\begin{itemize}
\raggedright
  \item \emph{Say it:} \emph{keitai}
  \item \emph{Say it:} \emph{keitai}, once more
  \item \emph{Say it:} \emph{keitai o uchi ni wasuremashita} — I left my mobile at home.
  \item \emph{Recall:} say the Japanese for to join, then the Japanese for a message (passed on), then say \emph{keitai} again
\end{itemize}

\begin{tcolorbox}[breakable,colback=teal!4,colframe=teal!35!black,title={Before you move on}]
What does \ja{けいたい} mean? (\textquotedblleft{}A mobile phone\textquotedblright{}.) Say it once more.
\end{tcolorbox}
\section[\texorpdfstring{\ja{かけなおす} (kakenaosu)}{かけなおす (kakenaosu)}]{\ja{かけなおす} (kakenaosu) — to call back}

\label{lesson:JA-C307-kakenaosu}



\begin{quote}
Before the new one: say the Japanese for a message (passed on), then the Japanese for a mobile phone.
\end{quote}

\subsection*{You'll want to know: \ja{かけなおす}}
\textbf{\ja{かけなおす}} — \emph{kakenaosu} — \textquotedblleft{}to call back\textquotedblright{}.

\begin{quote}
\textbf{\ja{あとで} \ja{かけなおします。}} — \emph{ato de kakenaoshimasu} — \textquotedblleft{}I will call back later.\textquotedblright{}
\end{quote}

Polite \textbf{\ja{かけなおします}}, \emph{kakenaoshimasu}; a godan verb. Past (polite): \textbf{\ja{ははが} \ja{でんわを} \ja{かけなおしました。}}, \emph{haha ga denwa o kakenaoshimashita}.

\begin{cousinweb}
\textbf{\ja{かけなおす}} (\emph{kakenaosu}) is a native Japanese word (wago), inherited rather than borrowed from Chinese, built from \emph{kake}, \textquotedblleft{}calling\textquotedblright{} (from \emph{kakeru}), and \emph{naosu}, \textquotedblleft{}to do again\textquotedblright{} (JMdict, KANJIDIC2).
\end{cousinweb}

\begin{grammarlens}[title={What you've built}]
One more word for telling what happened.
\end{grammarlens}

\subsection*{Your turn}
\begin{itemize}
\raggedright
  \item \emph{Say it:} \emph{kakenaosu}
  \item \emph{Say it:} \emph{kakenaosu}, once more
  \item \emph{Say it:} \emph{ato de kakenaoshimasu} — I will call back later.
  \item \emph{Recall:} say the Japanese for a message (passed on), then the Japanese for a mobile phone, then say \emph{kakenaosu} again
\end{itemize}

\begin{tcolorbox}[breakable,colback=teal!4,colframe=teal!35!black,title={Before you move on}]
What does \ja{かけなおす} mean? (\textquotedblleft{}To call back\textquotedblright{}.) Say it once more.
\end{tcolorbox}
